\chapter{Case Study: Real-Time Fraud Detection for Card Payments}
\label{cha:case-study}

The previous chapters described the Function Registry and resolution cascade in isolation from
any particular domain (Chapter~\ref{ch:proposed_solution}), traced their implementation
(Chapter~\ref{cha:impl}), and quantified their overhead (Chapter~\ref{cha:evaluation}). This
chapter grounds those mechanisms in a single, illustrative scenario drawn from retail banking:
a card-payment authorization workflow that screens every transaction for fraud before it is
settled. The scenario is not a validated production deployment; it is a realistic composition of
the \texttt{FraudPipeline} example already introduced in Section~\ref{sec:developer-workflow},
extended to internalize both scoring functions and to make an authorization decision, used here
to walk through the integration end to end and to discuss where it fits --- and where it does not
yet fit --- an actual banking environment.

The chapter is organised as follows. Section~\ref{sec:case-scenario} sets out the business
requirements a payment-authorization workflow imposes and why they line up with the problem this
work addresses. Section~\ref{sec:case-design} defines the workflow in the OmniFlow DSL.
Section~\ref{sec:case-deployment} walks the resolution cascade through two deployments of that
same workflow --- first to AWS, then to GCP --- to make the provider-symmetry claim concrete.
Section~\ref{sec:case-runtime} follows a transaction through the deployed workflow at runtime, and
Section~\ref{sec:case-discussion} discusses the scenario's relevance and its limits, connecting
back to the gaps already identified in Chapter~\ref{cha:conclusions}.

\section{Scenario and Requirements}
\label{sec:case-scenario}

Consider a bank that authorizes card payments through a workflow rather than a monolithic
service: each transaction is screened by two independent scoring functions before the bank
decides whether to settle it. \texttt{fraud-check} is a rules/ML-based service that flags
transactions matching known fraud patterns; \texttt{risk-score} computes a numeric risk score from
transaction and account features. Both are internal functions owned by the bank's fraud team,
deployed and iterated on independently of the orchestration logic. A third, external call settles
approved transactions against the bank's core banking system, which is a distinct system with its
own release cycle and is not managed through the registry.

Three requirements are characteristic of this setting and motivate the integration described in
this dissertation:

\begin{itemize}
    \item \textbf{Frequent, independent model updates.} Fraud-detection models are retrained and
    redeployed far more often than the orchestration logic that calls them --- sometimes daily, in
    response to newly observed fraud patterns. If the workflow definition embedded a concrete
    endpoint for \texttt{fraud-check}, every redeployment of the model would risk requiring a
    matching workflow redeployment, exactly the coupling problem motivating the Function Registry
    (Section~\ref{subsec:endpoint-resolution}).
    \item \textbf{Multi-cloud posture.} Operational-resilience regulation applicable to EU banks,
    notably the Digital Operational Resilience Act (DORA), pushes institutions to avoid
    concentration risk on a single cloud provider for critical functions. A fraud-screening
    workflow is a plausible candidate for a documented ability to run on a second provider, which
    requires the orchestration and the functions it calls to be portable, not just the
    infrastructure underneath them.
    \item \textbf{Auditability of bindings.} A bank's change-management process typically requires
    a record of which concrete function version and endpoint served a given deployment, which the
    registry's \texttt{updatedAt}-stamped entries provide in embryonic form
    (Section~\ref{subsec:registry-evolution}), even though a production deployment would need the
    stronger guarantees discussed in Section~\ref{sec:case-discussion}.
\end{itemize}

\section{Workflow Design}
\label{sec:case-design}

Listing~\ref{lst:case-payment-workflow} defines \texttt{PaymentAuthorization}. The first two
steps call \texttt{fraud-check} and \texttt{risk-score} through \texttt{internalFunction}, so
their endpoints are resolved through the registry at deployment time rather than hard-coded, as
described in Section~\ref{subsec:endpoint-resolution}. Unlike the external \texttt{risk-score}
variant shown in Listing~\ref{lst:wf-external-part}, both scoring calls here are internal and
therefore registry-managed --- a distinction to a partner-owned endpoint the bank happens to also
call. Their results are captured in \texttt{fraudResult} and \texttt{riskResult}. A third step,
\texttt{decision}, is a conditional (\texttt{switch}) step: it jumps to
\texttt{decline-transaction} if the fraud check flagged the transaction or the risk score exceeds
a threshold, and otherwise falls through to the \texttt{default}, \texttt{approve-transaction}.
The approve path issues an external call to the bank's own settlement endpoint, declared with a
literal \texttt{host}/\texttt{path} and therefore bypassing the registry entirely, exactly like
\texttt{external-risk-api} in Listing~\ref{lst:wf-external-part}: the settlement system is not a
function this integration deploys or resolves.

\begin{lstlisting}[language=Kotlin,caption={PaymentAuthorization workflow: two internal scoring calls feeding a decision.},label={lst:case-payment-workflow},basicstyle=\ttfamily\small,breaklines=false]
val PaymentAuthorization = workflow {
  name("PaymentAuthorization")
  description("Real-time fraud screening for card payment authorization")
  input("transaction")
  steps(
    step {
      name("fraud-check")
      context(call {
        method(POST)
        internalFunction("fraud-check", "./deploy/fraud-check.json")
        body(variable("transaction"))
        result("fraudResult")
      })
    },
    step {
      name("risk-score")
      context(call {
        method(POST)
        internalFunction("risk-score", "./deploy/risk-score.json")
        body(variable("transaction"))
        result("riskResult")
      })
    },
    step {
      name("decision")
      context(switch {
        conditions(
          condition {
            match(variable("fraudResult.body.flagged") equalTo value(true))
            jump("decline-transaction")
          },
          condition {
            match(variable("riskResult.body.score") greaterThan value(75))
            jump("decline-transaction")
          }
        )
        default("approve-transaction")
      })
    },
    step {
      name("approve-transaction")
      context(call {
        method(POST)
        host("core-banking.internal.example")
        path("/v1/settlements")
        body(variable("transaction"))
        result("settlementResult")
      })
    },
    step {
      name("decline-transaction")
      context(assign {
        variables(variable("decision") equalTo value("declined"))
      })
    }
  )
  result("decision")
}
\end{lstlisting}

Nothing in Listing~\ref{lst:case-payment-workflow} names a cloud provider, a region or a concrete
endpoint for either scoring function: that information is supplied once, at deployment time,
through the QuickFaaS descriptors referenced by \texttt{internalFunction} and the choice of
provider deployer. This is what makes the same definition deployable to either provider, as
Section~\ref{sec:case-deployment} shows next.

\section{Deploying the Same Workflow to Two Providers}
\label{sec:case-deployment}

Suppose neither \texttt{fraud-check} nor \texttt{risk-score} has been deployed yet, and the bank
deploys \texttt{PaymentAuthorization} to AWS first. OmniFlow bootstraps an AWS registry by
listing Lambda functions across every enabled region, finds it empty for both functions, and the
cascade reaches Level~3 for each: QuickFaaS deploys \texttt{fraud-check} and \texttt{risk-score}
as Lambda functions from the descriptors referenced in the workflow, following the AWS deployment
sequence detailed in Figure~\ref{fig:aws-deploy-sequence} --- IAM role resolution, subprocess
invocation, \texttt{GetFunction} polling and Step Functions invoke permissions. Both ARNs are
written to the registry, and \texttt{PaymentAuthorization} is deployed as a state machine.

A week later, the fraud team retrains and redeploys \texttt{fraud-check} with an updated model,
using QuickFaaS directly against the same \texttt{functionRef}; the redeployment updates the
Lambda's code (\texttt{UpdateFunctionCode}) without changing its ARN, so the registry entry stays
valid. The next time \texttt{PaymentAuthorization} is deployed --- unchanged, from the developer's
point of view --- \texttt{fraud-check} resolves as a Level~1 registry hit and no redeployment is
triggered by the workflow deployment itself; the workflow definition never had to change to pick
up the new model.

Now suppose the bank wants a documented ability to run the same fraud-screening workflow on a
second provider, per the multi-cloud requirement in Section~\ref{sec:case-scenario}. Because
\texttt{PaymentAuthorization} names no provider, deploying it to GCP requires no change to
Listing~\ref{lst:case-payment-workflow} --- only different QuickFaaS descriptors (targeting Cloud
Run instead of Lambda) and the GCP deployer. On this first GCP deployment, the GCP registry is
bootstrapped separately from the AWS one, both scoring functions miss at Level~1 and Level~2, and
Level~3 deploys them as Cloud Run services instead, following the ZIP-based deployment path
described in Section~\ref{subsec:quickfaas-aws-provider}'s GCP counterpart
(Section~\ref{sec:background_quickfaas}). The two registries stay independent, one per provider
and target account, so the same \texttt{functionRef} resolves to a Lambda ARN under one deployment
and to a Cloud Run URL under the other --- the uniform \texttt{url} field discussed in
Section~\ref{subsec:registry-evolution} is precisely what makes this substitution transparent to
the workflow definition.

\section{Runtime Invocation}
\label{sec:case-runtime}

Once deployed, a transaction flows through \texttt{PaymentAuthorization} differently depending on
which deployment served it, even though the DSL is identical. On the AWS deployment, Step
Functions invokes \texttt{fraud-check} and \texttt{risk-score} through the native
\texttt{arn:aws:states:::lambda:invoke} integration described in Section~\ref{sec:runtime-invocation}
and shown in Figure~\ref{fig:aws-runtime-invocation}, with the ARN fixed at deployment time; the
\texttt{decision} step is rendered as a native \texttt{Choice} state evaluating
\texttt{fraudResult}/\texttt{riskResult}, and \texttt{approve-transaction} is rendered as the
\texttt{apigateway:invoke}-style HTTP task, since its endpoint was declared with a literal
\texttt{host}/\texttt{path} rather than resolved through the registry. On the GCP deployment, the
same two scoring calls are rendered as HTTP calls to the resolved Cloud Run URLs, since GCP has no
Lambda-style native invoke equivalent in this integration. In both cases, the decision logic and
the settlement call are unaffected by which provider deployed the scoring functions: only the
\texttt{Resource}/\texttt{call} shape the renderer emits for the two internal steps differs, which
is exactly the boundary Section~\ref{sec:runtime-invocation} identified as provider-specific.

\section{Discussion}
\label{sec:case-discussion}

This scenario exercises the integration's central claim in a setting where it plausibly matters:
a bank can retrain and redeploy a fraud model on its own cadence without touching the
orchestration that calls it, and can retarget the same workflow definition at a second cloud
provider by changing deployment descriptors rather than the workflow itself. Both properties speak
directly to the two motivating forces in Section~\ref{sec:case-scenario} --- independent model
lifecycles and a documented multi-provider posture --- without requiring the fraud or platform
teams to hand-maintain endpoint bindings across redeployments.

The scenario equally exposes why this remains illustrative rather than production-ready. The
gaps already identified in Section~\ref{sec:concl-critical} are not hypothetical here: a registry
holding the live endpoint of a bank's fraud-screening function, writable by any deployment process
with filesystem access and carrying no integrity guarantee, is not a control a bank's
change-management or security function would accept as-is, and the AWS path's broad, automatically
granted IAM role is similarly at odds with a least-privilege environment handling payment data.
The applicability discussion in Section~\ref{sec:concl-applicability} and the corresponding future
work in Section~\ref{sec:concl-future} --- a hardened, access-controlled registry backend with an
attributable write history, and least-privilege IAM under change control --- describe exactly the
hardening this scenario would require before the workflow in
Listing~\ref{lst:case-payment-workflow} could authorize real transactions.
